\chapter{Apresentação}

O presente Projeto Pedagógico do Curso (PPC) de \ppccurso da Universidade Federal de Uberlândia foi elaborado para a implantação do currículo \ppcversao, com o propósito de consolidar uma formação tecnológica aplicada, crítica e socialmente engajada, alinhada às Diretrizes Curriculares Nacionais Gerais para a Educação Profissional e Tecnológica~\cite{cne_cp_1_2021}, ao Catálogo Nacional de Cursos Superiores de Tecnologia (CNCST, 4ª edição)~\cite{mec_cncst_2024} e às orientações institucionais da UFU~\cite{ufu_guia_ppc_2021}.

Este PPC responde à crescente demanda por profissionais capazes de projetar, implementar, operar e manter sistemas automatizados industriais em um cenário de transformação digital da indústria — a chamada Indústria 4.0 —, marcado pela convergência entre eletrônica, controle, redes industriais, robótica e inteligência artificial aplicada. O curso foi construído de forma coletiva pelo Núcleo Docente Estruturante (NDE), pelo Colegiado do Curso e por representantes da comunidade acadêmica, com a convicção de que a formação de um(a) Tecnólogo(a) em Automação Industrial deve ser eminentemente aplicada, profissionalizante e orientada por competências diretamente ligadas ao exercício profissional.

A denominação regulatória do curso — \ppccurso — está em plena aderência ao CNCST, eixo tecnológico Controle e Processos Industriais, área tecnológica Eletrônica e Automação (código CINE Brasil 0714A01)~\cite{inep_cine_brasil_622_2025}, em conformidade com a Resolução CNE/CP nº 1/2021~\cite{cne_cp_1_2021} e com a estrutura de eixos/áreas tecnológicas do CNCST fixada pela Resolução CNE/CP nº 2/2024~\cite{cne_cp_2_2024}. A ênfase curricular do curso — \ppcenfasecurricular\ — e o nome de divulgação institucional — \ppccursomercadologico\ — expressam a identidade pedagógica proposta neste PPC sem descaracterizar o curso perante o MEC/INEP (ver detalhamento em \autoref{chp:identificacao} e a justificativa da denominação em \autoref{chp:justificativa}).

Este PPC reflete um currículo estruturado em núcleos de formação básica, tecnológica e profissional, humanística, de integração curricular/extensão e de formação complementar, combinando disciplinas de fundamentos científicos, tecnologias centrais de automação industrial (eletricidade, eletrônica, instrumentação, controle, CLP/SCADA, redes industriais, acionamentos, hidráulica e pneumática) e um eixo diferenciador em robótica, sistemas inteligentes e inteligência artificial aplicada à automação. Mais do que uma exigência normativa, este documento representa a carta de intenções de um curso que se projeta para atender à demanda do setor produtivo por profissionais de formação tecnológica aplicada, sem perder de vista o rigor técnico-científico e a responsabilidade social de sua atuação.

O \ppccurso compõe uma proposta de Área Básica de Ingresso (ABI) com o curso de Engenharia de Automação, Robótica e Inteligência Artificial: os cinco primeiros períodos das matrizes são comuns e a proposta prevê a escolha da trajetória após essa etapa. A organização preserva a duração e a identidade próprias de cada curso — o \ppccurso integraliza-se em \ppctempominimo, com um sexto período específico. O ingresso, a escolha e eventual continuidade em uma segunda graduação dependem de ato institucional específico, conforme o \autoref{chp:abi}.
